\chapter*{\lblIntroduction}
\phantomsection
\addcontentsline{toc}{chapter}{\lblIntroduction}
\markboth{\lblIntroduction}{\lblIntroduction}

\section*{Dla kogo jest ta książka}

Wdrażałeś serwisy w C\#. Wiesz, po co jest pula połączeń, dlaczego
idempotencja ma znaczenie przy ponowionym żądaniu, co robi token anulowania
i mniej więcej na co spojrzałbyś najpierw, gdy serwis dławi się pod
obciążeniem. Poproszono cię \dash{} ogłoszenie o pracę, zespół, do którego
dołączyłeś, fakt, że ciekawe narzędzia w twojej dziedzinie są teraz pisane w
nim \dash{} żebyś robił to samo w Pythonie. I odkryłeś, że dwa tygodnie
wystarczą, żeby być produktywnym, i wcale nie wystarczą, żeby przestać pisać
Pythona, którego recenzent po cichu przepisuje.

Ta książka jest na te dwa tygodnie i na rok po nich. Zakłada, że jesteś
doświadczonym inżynierem, i nie zakłada, że znasz Pythona. Jest słownikiem
tłumaczeniowym z C\# na Pythona i podręcznikiem polowym wdrażania Pythona na
produkcję, napisanym przez kogoś, kto robi to zawodowo.

\section*{Czym ta książka nie jest}

\begin{itemize}[leftmargin=1.4em]
  \item Nie jest kursem Pythona od zera. Nigdy nie tłumaczy, czym jest pętla,
        a swoje strony wydaje na miejsca, w których dekada C\# daje Pythona,
        który działa i źle się czyta.
  \item Nie jest katalogiem bibliotek. Wybiera jedno narzędzie do każdego
        zadania \dash{} \pkg{uv}, \pkg{ruff}, \pkg{pydantic}, \pkg{FastAPI},
        \pkg{SQLAlchemy}, \pkg{pytest} \dash{} przypina wersję i mówi
        dlaczego. Alternatywy dostają jedno zdanie.
  \item Nie jest kolejną książką o AI. Część~V jest krótka i zatrzymuje się
        na jednym wywołaniu modelu i jednym ustrukturyzowanym wyniku.
  \item Nie jest książką o LangChain. Tamta książka, \emph{LangChain, LangGraph
        and Async Python}, jest już napisana dla inżyniera przychodzącego z
        .NET, a ta jest jej tomem zerowym: język, narzędzia, idiomy i nawyki
        produkcyjne, które tamta zakłada. Frameworki agentowe zostają tam.
        Microsoft Agent Framework zostaje w trzecim tomie. Razem te trzy mają
        być czytane jako jeden kurs i żadna nie powtarza sekcji innej.
\end{itemize}

\mermaidfig{book-map}{Gdzie jest ta książka: tom zerowy trylogii, z tym samym czytelnikiem na obu stosach}{trylogia}

\section*{Dlaczego teraz}

Trzy rzeczy zmieniły się w ciągu dwóch lat przed napisaniem tej książki i
każda odebrała jedną wymówkę.

Python~\pyver{} uczynił wersję free-threaded oficjalnie wspieranym sposobem
uruchamiania interpretera, a nie eksperymentem, co znaczy, że jedno zdanie,
które każdy inżynier .NET zna o Pythonie \dash{} \enquote{ma globalną blokadę
interpretera} \dash{} potrzebuje teraz drugiego zdania, a
Rozdział~\ref{ch:runtime} je dostarcza i mierzy.

Narzędzia się zbiegły. Przez większość historii Pythona uczciwą odpowiedzią na
\enquote{co jest odpowiednikiem \code{dotnet restore}} było wzruszenie ramion i
lista; odpowiedzią jest teraz \pkg{uv}, a Rozdział~\ref{ch:packaging} jest
napisany tak, jakby to było rozstrzygnięte, bo jest.

A historia typowania dojrzała do punktu, w którym instynkty inżyniera C\# są
atutem, a nie obciążeniem, pod warunkiem że nauczy się dokładnie, gdzie
checker się zatrzymuje. To Rozdział~\ref{ch:typing} i to dla niego ta książka
została zaczęta.

\section*{Jak zorganizowana jest książka}

Pięć części, czternaście rozdziałów, każdy poniżej trzech tysięcy słów.

\textbf{Część~I \dash{} \lblPartI} (Rozdziały~\ref{ch:runtime}
do~\ref{ch:typing}) to miejsce, w którym inżynierowie .NET przewracają się
pierwsi: środowisko uruchomieniowe, środowisko projektu, pakowanie, typowanie.
Nic z tego nie jest trudne i wszystko jest inne, a zrobienie tego źle daje
półroczny projekt, którego nikt nie umie zainstalować.

\textbf{Część~II \dash{} \lblPartII} (Rozdziały~\ref{ch:objects}
do~\ref{ch:imports}) to słownik właściwy: obiekty, funkcje, błędy, importy.
Każdy rozdział to tabela mapowania i powody, dla których mapowanie nie jest
jeden do jednego.

\textbf{Część~III \dash{} \lblPartIII} (Rozdział~\ref{ch:asyncio}) to jeden
rozdział, celowo, bo książka o LangChain spędza już dwa na pętli zdarzeń dla
obciążenia agentowego. Ten jest tłumaczeniem z Task Parallel Library i jest
rozdziałem, o którym inżynierowie .NET myślą, że go nie potrzebują.

\textbf{Część~IV \dash{} \lblPartIV} (Rozdziały~\ref{ch:services}
do~\ref{ch:observability}) to produkcja: serwis, baza danych, zestaw testów i
powierzchnia operacyjna \dash{} logi, ślady, kontenery, audyty \dash{} każde
zmapowane z odpowiedników w ASP.NET Core i EF Core.

\textbf{Część~V \dash{} \lblPartV} (Rozdziały~\ref{ch:aitoolkit}
i~\ref{ch:traceassert}) to zestaw inżyniera AI \dash{} trzy rzeczy, z których
zbudowany jest każdy SDK w tej dziedzinie \dash{} i ukończenie projektu
przewodniego.

Rozdziały są na tyle niezależne, że po części~I można je czytać w dowolnej
kolejności, z dwoma wyjątkami: Rozdział~\ref{ch:testing} zakłada
Rozdział~\ref{ch:functions}, a Rozdział~\ref{ch:traceassert} zakłada
wszystko.

\section*{Projekt przewodni}

Książka buduje przez całą swoją długość jedną rzecz:
\textbf{\code{trace-assert}}, port na Pythona pierwszej warstwy
\code{agent-eval-bench} \dash{} deterministycznych asercji nad śladem
wykonania agenta, napisanych jako testy \pkg{pytest}. \emph{Czy agent wywołał
narzędzie, które miał wywołać? W tej kolejności, w której miał? Z argumentami,
które przechodzą walidację? Nie wywołując niczego, czego mu zabroniono?} To
pytania, na które test odpowie bez modelu, bez budżetu i bez osądu, i są
najtańszą ewaluacją, jaką agent może mieć.

Przychodzi w trzech etapach. Rozdział~\ref{ch:testing} kładzie szkielet
\dash{} model śladu, fixture i dwie pierwsze asercje.
Rozdział~\ref{ch:aitoolkit} zapisuje prawdziwe wywołanie modelu jako
zdarzenia śladu. Rozdział~\ref{ch:traceassert} domyka zestaw asercji, pakuje
go i kończy tabelą porównującą, ile ten sam instrument kosztował w .NET, w
TypeScripcie i w Pythonie.

Projekt nie jest zabawką wybraną do nauczania. Jest rzeczą, której autor
potrzebował, a czytelnik kończy książkę z pakietem, wpisem na blog i
odpowiedzią na rozmowie rekrutacyjnej, co jest uczciwym powodem, żeby w ogóle
wybierać jakiś projekt.

\section*{Zasady, których ta książka się trzyma}

\begin{note}
\textbf{Każdy listing jest plikiem, a każdy plik jest uruchamiany.} Kod
mieszka w katalogu \code{code/} jako osobny projekt, zablokowany na wersjach
ze strony tytułowej, a build repozytorium uruchamia każdy listing i każdy test
ćwiczenia przy każdej zmianie. Listing, który przez to nie przeszedł, mówi o
tym w ramce.
\end{note}

\begin{note}
\textbf{Pomiary ponad twierdzenia.} Twierdzenie o tym, co jest szybsze,
tańsze albo bezpieczniejsze, przychodzi ze skryptem i liczbą albo jest
oznaczone jako osąd. Osiem pomiarów jest wyspecyfikowanych w rozdziałach,
wszystkie darmowe i kończące się na laptopie; dopóki ten za twierdzeniem nie
został uruchomiony, twierdzenie o tym mówi.
\end{note}

\begin{note}
\textbf{Twarde ramy.} Czternaście rozdziałów, żaden powyżej trzech tysięcy
słów, każdy publikowalny osobno, w trzech wydaniach. Ramy istnieją, bo inne
projekty autora nauczyły go, czym staje się książka bez granic, a build je
egzekwuje, zamiast ufać, że ktokolwiek będzie pamiętał.
\end{note}
